% sec-03: Document 03, SB 964, model transparency and clinical data protection.
\docpart{SB 964 Data Protection}{SB 964 - California Oncology Model Transparency
and Clinical Data Protection Act of 2026}

\section{Findings and Declarations}

\draftnote{Six findings. Anchor them to the California instruments already on
the books rather than inventing a policy rationale: AB 3030 on artificial
intelligence in health care services, AB 489 on deceptive terms, SB 1120 on
utilization review, and AB 2013 on training data transparency. State what this
bill adds that those four do not reach, which is the clinical trial setting and
the retention of prompts and completions as source records.}

\section{Definitions}

\draftnote{Define model, model version, prompt, completion, hash-chained audit
trail, provenance record, and secondary use. Take the audit-trail construction
from \mvfile{national-platform/} and keep the definition identical to the one
used in document 07 \sect{}5.}

\section{Model Transparency Obligations}

\begin{mvtable}
\begin{tabularx}{\textwidth}{>{\raggedright\arraybackslash}p{3.6cm} >{\raggedright\arraybackslash}p{2.6cm} >{\raggedright\arraybackslash}X}
\headrow \thead{Obligation} & \thead{Frequency} & \thead{Record that discharges it}\\
\midrule
Training data provenance disclosure & to be fixed & to be fixed \\
Model version register & to be fixed & to be fixed \\
\end{tabularx}
\end{mvtable}
\tabcapl{tab:sb964obl}{Transparency obligations. Stage 2 completes the table
against the model governance section of document 07.}

\draftnote{Write five obligations. Each must be discharged by a record that
exists, not by a policy statement. Cross-check the register requirement against
\mvfile{funding/move-in/draft-move-in/sections/sec-07-fda-compliance-guide.tex}
so the state register and the federal change-control record are the same
artifact rather than two.}

\section{Clinical Data Protection}

\draftnote{De-identification standard, the prohibition on secondary use without
specific consent, and the relationship to 45 CFR Part 164. State which of the
two standards governs where they differ, and do not leave the reader to work it
out.}

\section{Audit Trail and Retention}

\draftnote{Retention periods aligned to 21 CFR 312.62 and to the archive
requirement in document 13 \sect{}6. Where the periods differ, the longer
governs; say so.}

\section{Enforcement and Breach Notification}

\draftnote{Notification clock, the department's remedies, and the interaction
with existing California breach law. Keep the deficiency classes identical to
documents 01 and 06.}
